\section{Account for the visible and numerical loss}\label{rgb-loss}
A five-bit field has 32 possible values and a six-bit field has 64. Green therefore
has twice as many retained code levels as either other component. The name RGB565
describes that allocation; it does not prescribe the conversion rule or a perceptual
error model. Here all channels use the explicit truncation and expansion formulas.

\Needspace{18\baselineskip}
\begin{table}[htbp]
\centering
\caption{Individual channel examples. Each pair uses its channel's specified quantizer and expansion rule.}
\tagpdfsetup{table/header-rows={1}}
\begin{tabular}{lrrr}
\toprule
Channel & Input & Retained field & Expanded output\\
\midrule
Red & 7 & 0 & 0\\
Red & 8 & 1 & 8\\
Red & 127 & 15 & 123\\
Red & 128 & 16 & 132\\
Green & 63 & 15 & 60\\
Green & 64 & 16 & 65\\
Blue & 247 & 30 & 247\\
Blue & 255 & 31 & 255\\
\bottomrule
\end{tabular}
\end{table}

Packing discards low bits, but the expanded value need not be smaller than the
original. Red 128 becomes 132 under bit repetition. Writing a five-bit input as
$8q+u$, where $0\leq u\leq7$, gives expansion error
$\lfloor q/4\rfloor-u$. It lies between minus seven and seven. For green, the
corresponding error lies between minus three and three. These bounds concern each
component's integer code value, not a distance in a perceptual color space.

The supplied preview places the original 32-pixel row above the expanded row.
Small numerical changes may be difficult to distinguish at a particular zoom or
on a particular display. Inspect bytes and computed errors as well as the image.
A matching screenshot cannot establish that every possible component value is correct.
